\section{Where the Knowledge Comes From}
\label{sec:authoring}

If the contract's value is the knowledge it states, its value is its
author's. Section~\ref{sec:authoring-paths} describes three ways to author one:
interactively, a model drafting under the author's direction; compiled, entry by entry, from the documentation; and authored
whole by a model that reads all of the documentation and explores the
data. This section compares them on both benchmarks.

\subsection{A first compile is worth the manual}
\label{sec:compiled}

The compiler reads a manual or KB entry by entry and turns each into a
contract definition, checking its SQL against the database read-only and
refusing what it cannot state; on LiveSQLBench every one of the 1{,}369
KB definitions is in the compiled contracts verbatim. Compiled knowledge,
delivered either way, does what the manual does. On LiveSQLBench's 302
primary tasks \mLuna{} scores 35.5\% with the KB as a manual, 36.8\% with
the manual plus every compiled metric's SQL (\armMsql{}), and 35.9\% with
the compiled contract behind the governed tools; contrary to its
registered prediction (3), the contract does not beat \armMsql{} (22 tasks
to 25). \mQwen{} gives
37.3\%, 37.7\% and 37.7\%. On DABStep the same compiler (Sonnet~5) yields
45.6\% on \mLuna{} over all tasks, against the manual's 45.3\% in an
earlier run (not paired, so not tested) and the original contract's
80.2\% in the same run (196 tasks to 18). On both benchmarks the
agent reads what it is given, looking a metric up in 86 to 93\% of
compiled-contract runs on LiveSQLBench, so the null is not a contract left
unread. A blind model review of the 355 compiled metrics that primary
tasks depend on, with the database but no task, found 247 faithful, 60
wrong in a way that changes results (13 wrong without that, 35 unclear),
most often a \texttt{NULL} silently counted or a unit misread.

\subsection{A better compile recovers what the data settles}
\label{sec:better-compile}

On DABStep the compiler is the limit, not the method. A later run put four
contracts side by side on \mLuna{}. With Opus~5.5 as the compiler and a contract format that can declare a multi-table metric, the
fee formula compiles, and the compiled contract reaches 76.6\% against the
Sonnet compile's 47.3\% (172 tasks to 15) and the original contract's
80.0\% in the same run. Opus without joins compiles a one-table
contract that scores 47.4\%; Sonnet with joins still refused the formula,
and was not run as an agent arm. The best compile still trails the original
contract (58 tasks to 22), by 3.4 points. It works entry by entry from the
manual, but its checks against the data found what the manual leaves out:
its notes record that the data stores ``all'' as an empty list and that
several fee rules can match one payment, and its fee definition returns
one row per matching rule. \textbf{A compiler that checks its definitions
against the data adds what the data settles}, here the wildcard of
Section~\ref{sec:delivery} and the overlap of rules behind the second fact,
though its note leaves open how overlapping fees combine.

On LiveSQLBench a better author made no detectable difference. The LSB
author experiment gave three databases (55 primary tasks) two new contracts: an Opus compile with
declared joins, and one written whole by an Opus author session that read
the entire KB and explored the data read-only, without seeing any task.
Against the Sonnet compile, neither differs: the Opus compile 5 tasks to
3 ($+1.8$ points, interval $-3$ to $+7$), the author 3 to 4 ($-1.8$, $-9$
to $+6$; intervals over tasks, since three databases are too few to
cluster), which on these tasks excludes a gain of more than about 7 points;
between any two contract arms only 7 or 8 of the 55 tasks score
differently. Section~\ref{sec:lsb-knowledge} says why. Most of the hard
knowledge there is not hard to write down; the KB does not settle it, so
no author can recover the gold's choice.

The LSB oracle experiment tested that directly, with oracle knowledge. Each database's
primary tasks were split in half before anything was written. From the
gold of the 27 training tasks only, 35 corrections were written as general
facts about KB entries (a scale, a threshold, how a missing value counts,
a grain), never a question's filter or output form, and applied both to
the manual and to the Opus contract. They lift the training tasks by
about 21 points in both deliveries (7 tasks to 0 each). On the 28 held-out
tasks they lift the manual by 2 tasks to 0 ($p{=}0.5$; $+3.6$ points,
interval 0 to $+9.5$) and the contract not at all (0 tasks to 2), which
rules out the training tasks' 21 points, although half of those tasks use an entry a correction names. A
choice learned from one task's gold does not carry to another's through
the same entry; two training golds already used one entry differently. On
LiveSQLBench the undetermined knowledge is settled task by task in the
gold, and a contract, which can only hold what holds across questions,
has nothing to hold there.

\subsection{An author with the data recovers what the data settles}
\label{sec:gaps}

\begin{table}[t]
\centering
\footnotesize
\setlength{\tabcolsep}{2pt}
\begin{tabular}{@{}lrrrrr@{}}
\toprule
& \textsc{gapped} & \textsc{full} & \textsc{authored} & \textsc{noted} & \shortstack{\textsc{authored}\\\textsc{text}} \\
\midrule
data-informed (41) & 47.7 & 97.6 & 74.5 & 75.3 & 74.0 \\
\quad settled (20) & 59.4 & 98.3 & 99.4 & 100.0 & 98.9 \\
\quad split (21) & 36.5 & 96.8 & 50.8 & 51.9 & 50.3 \\
convention-only (9) & 0.0 & 97.5 & 0.0 & 0.0 & 0.0 \\
\bottomrule
\end{tabular}
\caption{The gaps experiment, accuracy (\%) pooled over \mLuna{},
\mQwen{} and \mSol{}, $k{=}3$; the 41 data-informed questions come from
five gaps (\textsc{authored} over \textsc{gapped}: 4 gaps to 0, one tied). Rows: questions by what the data reveals
about their gap (defined in the text); convention-only questions need the
precedence rule. \textsc{gapped}: the KB with gaps opened;
\textsc{full}: the complete KB; \textsc{authored}: the contract an author
wrote from the gapped KB with the data open; \textsc{noted}: the gapped
KB plus that author's resolutions as notes; \textsc{authored text}: the authored files in the prompt.}
\label{tab:gaps}
\end{table}

Real documentation leaves things open. The gaps experiment opened gaps in
the complete KBs of the constructed questions and asked whether an author with the data
recovers them, and whether the contract's gain is then the author's
knowledge. The gaps were a wildcard stored as a sentinel (in all three
databases), a quantity stored in an encoding the KB does not state (a
score one higher, a fraction where the KB implies a percentage), and, in
one database, the precedence among matching rules removed. Blind model
raters labelled each gap by what the data reveals: \emph{settled},
\emph{split}, or \emph{convention-only}. The author wrote three
contracts per database in independent sessions with the data open, and
one without the data as a control.

Recovery was scored before any agent ran. Over five gaps and three
sessions, the sessions with data recovered 12 of 15: settled 6 of 6, split
6 of 9, convention-only 0 of 3. The session without data recovered none of
five. Every session, with data or without,
resolved the precedence gap as the lowest rule id; the gold is the most
specific rule. The agents then show the recovered knowledge doing the
work (Table~\ref{tab:gaps}). Five arms ran on the three models: the KB
with the gaps opened (\textsc{gapped}), the complete KB (\textsc{full}),
the contract the first session wrote (\textsc{authored}), the gapped KB
plus that session's resolutions written as notes (\textsc{noted}), and the
authored contract's files in the prompt (\textsc{authored text}). The
authored contract beats the gapped KB by 26.8 points on the 41
data-informed questions, as registered in (35): 28 tasks to 1, but those
questions come from five gaps, and over gaps the gain holds on four with
one tied (one-sided $p{=}0.06$, the smallest five clusters allow being
0.03). By model it is $+50$ points on \mLuna{}, $+26$ on \mSol{} and $+4$ on
\mQwen{} (6 tasks to 2, not significant). The \emph{knowledge share}, the
part of that gain \textsc{noted} reproduces, is 1.03; its 95\% interval,
0.98 to 1.35, is bootstrapped over only five gap clusters and is a rough
guide. The three arms that hold the author's knowledge tie on
every label, whatever their form. And the author's misses travel with the
contract: on the encoding the data only suggested and on the precedence
convention it cannot reveal, every arm built from that author scores 0,
and the complete KB 97 to 98\%.

\subsection{What decides when a contract pays}
\label{sec:account}

Three accounts were tested and dropped along the way. Reuse is not it: a
LiveSQLBench database's most reused KB entry serves a median 32\% of its
tasks against DABStep's 44\%, and the contract's gain does not rise with
reuse on either model. The type of knowledge is not it: a rule-matching
procedure, stated completely, gets through prose as well as through a
contract (Section~\ref{sec:procedures}). And delivery, once the facts are equal,
adds a few points on the rule-matching tasks (Section~\ref{sec:delivery}), not the
20 to 35 at issue. What survives is completeness and
determinability. \textbf{A contract pays where it states something true
that the prose leaves out. An author can supply what the documentation or
the data determines, and a convention neither determines only when the
gold happens to use the author's default.} DABStep shows both. The
empty-list wildcard is settled by the data. That the fees of overlapping
rules add up is a convention, and the original contract and the Opus
compile both wrote the ordinary join-and-sum, which is the gold's; so the
original contract leads the first compile by 35 points over all tasks on
\mLuna{}. The gaps experiment's precedence rule is a convention too, but
its gold is not the default, and every author missed it. LiveSQLBench's
hard knowledge is settled task by task in the gold, so no author we tried
gained detectably.
